% !TEX root = ../main.tex
\chapter{결론과 앞으로 할 일}
\label{ch:conclusion}

\ref{ch:discussion}장에서는 발견을 풀이하고, 이 연구가 목표(O)를 이루었는지 따져 보았어요.

\dissertationSubheading[sec:summary-findings]{연구 질문에 대한 답}

이 논문은 결합 페이지랭크(C-PR, 두 그래프를 함께 걷는 점수)를 만들었어요. C-PR은 ERC-20(토큰 공통 규칙) 허락(allowance, 토큰을 얼마까지 대신 써도 된다는 허락) 화살표(엣지)와 송금(transfer) 화살표 위를 걷는 하나의 걷기예요. 이때 두 종류의 화살표는 주소마다 정해 둔 무게 $\lambda$로 섞어요. 이 논문은 C-PR을 한 층 비교 기준들과 나란히 놓고 평가했어요. 하나는 엔도스랭크(EndorseRank)예요. 엔도스랭크는 마지막 허락(가장 최근 허락) 위에서 AWP(적응형 가중 페이지랭크)의 화살표 무게(가중치)를 쓰는 페이지랭크(PageRank)예요. 다른 하나는 AWP 자체예요 \parencite{do2023awp}. 또 C-PR의 빼 보기 실험인 S-PR(보증 씨앗 페이지랭크), 그리고 원시 차수(화살표를 그냥 센 수)와도 나란히 놓고 평가했어요. 평가에는 아비트럼 원(Arbitrum One)의 스마트 계약(블록체인 위에서 저절로 실행되는 프로그램)을 썼어요. 점수는 동결일(점수를 고정하는 날)마다 그 앞 열두 달의 기록으로 매겼어요. 홀드아웃(나중 기간을 떼어 두고 맞혀 보는 시험) 기간은 W0(2026년 4월부터 6월까지)과 W1(2026년 7월부터 9월까지)이에요. 다섯 토큰의 달러(USD) 금액, 코호트(묶음), 점수 기간은 등록 분석(미리 등록한 그대로 한 분석)의 결과를 안 뒤에 정했어요. 아래의 규칙 판정은 재분석이에요. 이 판정들은 R1을 빼면 등록 분석의 판정과 같아요(부록~\ref{sec:map-registered}).

RQ1은 엔도스랭크와 AWP가 한 표본의 계약들을 서로 다르게 줄 세우는지 물었어요. 엔도스랭크는 허락 화살표만 걷고, AWP는 송금 화살표만 걸어요. 그리고 RQ1은 각 점수가 같은 기간의 허락 통계와 송금 통계에 얼마나 가깝게 맞는지도 물었어요. 두 점수는 계약들을 서로 다르게 줄 세워요. 그리고 가설(H) H1은 뒷받침돼요. 다만 대부분은 만든 방식 때문에 저절로 그래요(\ref{sec:rank-divergence}절). 각 점수는 자기가 걷는 화살표를 따라가요(표~\ref{tab:alignment-hybrid}). 이것은 H2를 좁은 뜻으로만 뒷받침해요. 대리 지표(대신 재는 값)가 그래프와 같은 이벤트(계약이 남기는 기록)에서 나오기 때문이에요. 이 갈라짐이 결합 걷기의 바탕(전제)이에요.

RQ2는 각 점수를 계산하는 데 비용이 얼마나 드는지, 그리고 두 화살표 묶음을 모두 걸으면 비용이 얼마나 더 드는지 물었어요. 비용은 각 그래프의 크기와, 그 걷기에 필요한 반복(한 바퀴) 수에 따라 커져요. 그래서 가장 듬성듬성한(화살표가 적은) 그래프 위의 엔도스랭크가 가장 싼 점수예요. 그리고 결합 걷기들은 AWP보다 비용이 더 들어요(표~\ref{tab:benchmark-runtime}). H3은 뒷받침돼요. 다만 대부분은 만든 방식 때문에 저절로 그래요. 또 부분 표본(코호트에서 일부만 뽑은 표본)에서도 엔도스랭크는 모든 단계에서 계속 더 싸요. 하지만 시간 비율이 화살표 수 비율을 따라가는 단계는 하나도 없어요(표~\ref{tab:benchmark-scaling}).

RQ3은 $t_1$에 동결한 점수가, 스펜더 계약(허락을 받아 대신 쓰는 계약)이 그 뒤에 얻는 새 허락 쌍(새로 생긴 주인–스펜더 허락)과 새 송금자(새로 송금을 보내온 주소)를 예측하는지(미리 맞히는지) 물었어요. 그리고 결합 걷기를 포함해 어떤 페이지랭크든, 같은 종류 화살표의 원시 차수보다 이것들을 더 잘 예측하는지도 물었어요. 첫째 부분은 두 기간 모두에서 성립해요. 그리고 모든 스펜더를 놓고 보면, 점수가 어느 라벨(나중에 채점에 쓰는 정답 값)을 미리 내다보는지는 화살표 종류가 정해요(표~\ref{tab:holdout-spenders}와~\ref{tab:fresh-holdout}). 둘째 부분의 답은 모든 걷기에서 “아니요”예요. 이것은 라벨을 본 뒤에 계산한 대비(두 값의 차이)에서 나온 답이에요. 두 기간 모두에서, AWP는 나중의 새 송금자를 송금 입차수(송금을 보내온 서로 다른 주소 수)보다 못하게 줄 세우고, 엔도스랭크는 나중의 허락 쌍을 받은 허락 수(허락 입차수)보다 못하게 줄 세워요(표~\ref{tab:holdout-diff}과~\ref{tab:fresh-posthoc}). AWP의 활동 재시작(활발한 주소에서 더 자주 다시 출발하기)을 쓰는 엔도스랭크는 차수에 모자라는 정도가 더 작아요. 그리고 C-PR은 두 라벨 어느 것도 같은 종류 화살표의 차수보다 더 잘 줄 세우지 못해요(표~\ref{tab:cpr-contrasts}). 다시 출발 규칙은 걷기가 자기 화살표를 그냥 센 수에 얼마나 모자라는지를 바꿀 뿐, 모자라는지 아닌지는 바꾸지 않아요.

RQ4는 이 논문의 방법인 C-PR이, 자기의 두 층 가운데 어느 한 층만 걷는 것보다 계약을 더 잘 줄 세우는지 물었어요. 분석 계획이 정해 둔 결합 방식인 C-PR에서 답은 “아니요”예요. 등록 분석에서처럼, 고친 점수에서도 그래요(표~\ref{tab:registered-vs-revised}). 규칙~A에서 C-PR은 주된 부분의 허락 조건은 채우지만, 그 송금 조건과 보조 부분은 채우지 못해요(\ref{sec:coupled-results}절; 표~\ref{tab:holdout-diff}). 씨앗을 쓰는 빼 보기 실험은 새 송금자를 송금 걷기보다 조금 더 잘 줄 세우고, 새 허락 쌍은 허락 걷기보다 훨씬 못하게 줄 세워요. 이런 형태에서는, 허락 화살표가 주로 걷기가 그 화살표를 따라 움직일 때 도움이 되고, 씨앗으로서는 거의 도움이 되지 않아요.

RQ5는 허락 층을 더하면 송금만 걷는 걷기가 나아지는지 물었어요. 이번에도 답은 “아니요”예요. 규칙~B에서, 새 허락 쌍에서 송금 걷기보다 나았던 W0의 탐색적(미리 정하지 않고 살펴본) 이득은 W1에서 다시 나타났어요. 이 이득의 일부는 송금 걷기가 0점을 주는 스펜더들에서 왔을 수도 있어요. 하지만 새 송금자에서의 손실은 허용 폭(더 나빠도 괜찮다고 미리 정한 폭)을 넘었어요. 그리고 트레이더(거래하는 계약)의 청산 없이 닫은 비율(청산, 곧 손실이 너무 커져 거래가 강제로 끝나는 일 없이 거래를 닫은 비율)에서는, 허용 폭보다 큰 손실이 없다고 말할 수 없었어요. AWP와 견주면, C-PR은 새 허락 쌍을 더 잘 줄 세우고 새 송금자는 더 못하게 줄 세웠어요. 그리고 그 닫은 비율에서는 어느 쪽이 나은지 가려지지 않았어요(표~\ref{tab:fresh-contrasts}). 규칙과는 따로, 라벨을 본 뒤에 계산한 대비를 보면 이래요. C-PR은 두 기간 모두에서 두 스펜더 라벨을 엔도스랭크보다 더 잘 줄 세워요. 나중의 허락에서 C-PR이 AWP보다 앞서는 것은 주로 AWP가 0점을 주는 계약들에서 나와요. W1에서는 나머지 계약들 사이에서도 조금 앞서요. 그리고 C-PR은 두 기간 모두에서 나중의 허락을 받은 허락 수보다 못하게 줄 세우고, W1에서는 활동 재시작을 쓰는 엔도스랭크보다도 못하게 줄 세워요(표~\ref{tab:cpr-contrasts}).

규칙~B와 함께 등록했지만 연구 질문을 더하지는 않는 비교가 있어요(\ref{sub:finding-neutral}절). 이 비교의 결과 값은 나중에 청산을 피하는 일이에요. 이 결과 값은 두 종류의 화살표 어느 쪽으로도 만들지 않았고, 거래 상대의 위험과 관계가 있어요. 주된 기간인 W1에서, 동결일의 허락한 주인 수(서로 다른 주인 수)는 이 결과 값에서 송금 수보다 앞섰어요($P=+0.141$, 97.5\% 구간(여러 번 비교할 때 넓힌 신뢰구간) $[0.026,0.256]$; \ref{sub:neutral-label-results}절). W0에서는 양수였지만 가려지지 않았어요. 같은 조건으로 맞춘 페이지랭크 쌍은 어느 쪽이 나은지 가려지지 않았어요. 또 이 신호는 거래를 잘하는 것이 아니라 강제 청산을 피하는 것에 대한 신호예요. 그리고 표본이 작아요.

주장~C1--C6(\ref{sec:empirical-claims}절)은 이 발견들을 정식으로 적어 놓은 거예요.

주된 연구 질문은, 두 종류의 화살표를 한 번에 걷는 걷기가 이 계약들을 줄 세우는 데 무엇을 보태는지 물었어요. 라벨을 본 뒤에 계산한 대비에서 보면 이래요. 허락만 걷는 걷기와 견주면, 결합 걷기는 송금 구성개념(점수가 재려고 하는 생각)을 보태고, 나중의 스펜더 라벨을 모두 더 잘 줄 세워요. 송금만 걷는 걷기와 견주면, 결합 걷기는 맡긴 권한을 보태요. 이 덕분에 나중의 허락 쌍을 더 잘 줄 세우는데, 주로 송금 걷기가 줄 세우지 못하는 계약들에서 그래요. 그 대신 나중에 들어오는 흐름에서는 손해를 봐요. 미리 정한 규칙 아래에서는 결합 걷기가 어느 층도 이기지 못해요. 그리고 나중에 계산한 대비에서는, 결합 걷기를 포함해 이 연구의 어떤 걷기도 기간 밖에서 같은 종류 화살표의 원시 차수를 이기지 못했어요.

\dissertationSubheading[sec:contribution-knowledge]{지식에 보탠 새로운 기여}

\ref{ch:literature}장의 앞선 연구들에 비추어 읽으면, \ref{sec:contributions}절의 다섯 가지 기여(이 연구가 새로 보탠 것)는 다음을 밝혀요.

이 연구 전에는, 누구나 볼 수 있는 장부(블록체인) 위의 그래프 평판이 지갑을 그 지갑의 송금으로 줄 세웠어요 \parencite{do2023awp,do2023hodgerank}. 또는 위험이 거래 네트워크를 따라 어떻게 퍼지는지로 줄 세웠어요 \parencite{lin2025riskprop}. 또는 드러내 놓고 매긴 평가로 줄 세웠어요 \parencite{bellini2020trust,almasoud2020reputation}. 그리고 디파이(DeFi, 탈중앙 금융)에서 쓰이는 평판 도구는 자격 증명을 바탕으로 주소를 가진 사람에게 점수를 매겼어요 \parencite{humanpassportnddocs}. ERC-20 허락은 토큰을 쓸 권한을 일부러 맡겼다는 온체인(블록체인 위의) 기록이에요 \parencite{vogelsteller2015erc20}. 하지만 이 연구들 가운데 어느 것도 ERC-20 허락을 평판 순위의 화살표로 쓰지 않았어요. 허락을 돈을 낸 기록(송금)과 함께 한 걷기 안에 묶지도 않았어요. 그리고 계약 계정을 그 계정이 받은 허락으로 줄 세우지도 않았어요. 허락은 한 주체의 주소들을 묶는 데 쓰였어요 \parencite{victor2020clustering}. 그리고 보안 위험으로서 연구되었어요 \parencite{wang2022unlimited}. 또 온체인 평판 점수를 이렇게 평가한 연구도 없었어요. 한 코호트에서 각 구성개념의 점검과 견주고, 점수가 다듬는 원시 차수와도 견주는 평가예요. 이 평가는 신뢰구간(믿을 만한 범위), 차이 검사, 시간 홀드아웃을 쓰고, 미리 정한 판정 규칙을 따라요. 이 논문은 이 둘을 모두 내놓아요. 페이지랭크와 AWP의 무게 매기기는 앞선 연구에서 가져왔어요 \parencite{page1999pagerank,do2023awp}. 결합 걷기는 이동 규칙(다음에 어느 점으로 갈지 정하는 규칙)을 바꾸고, 결합 걷기의 허락 쪽 비교 기준인 엔도스랭크는 화살표만 바꿔요.

새로 얻은 지식은 다섯 부분으로 되어 있어요. 첫째, 계약의 온체인 평판 점수가 무엇을 재는지는 그 점수가 쓰는 화살표에 달려 있어요. 그리고 이것은 한 코호트에서, 불확실한 정도를 수로 나타내면서 점검할 수 있어요. 둘째, 허락 화살표를 송금과 묶는 걷기는 미리 정한 규칙 아래에서 자기의 두 층 어느 것도 이기지 못해요. 이렇게 묶는 것은 여러 층 네트워크 연구들이 제안한 한 걸음이었어요. 규칙과는 따로 보면, 이 걷기는 나중의 스펜더 라벨을 모두 허락만 걷는 걷기보다 더 잘 줄 세워요. 그리고 송금만 걷는 걷기에는 주로 한 가지 순서를 보태요. 송금 걷기가 동점(점수가 똑같은 것)으로 남겨 두는 계약들 사이의 순서예요. 그 대신 나중에 들어오는 흐름에서는 손해를 봐요. 셋째, 허락 화살표 하나만으로 만든 점수는 새로 계산하는 비용이 더 적어요. 이 점수는 송금에 바탕을 둔 점수와 다르고, 두 기간에서 스펜더가 나중에 받는 보증(믿고 맡김)을 줄 세워요. 그리고 송금 기록이 없는 스펜더 계약들에도 순서를 매겨요. 송금 걷기는 이런 계약들에 0점을 줘요. 나중의 보증에서 이 점수가 AWP보다 앞선다는 사후(라벨을 본 뒤의) 결과는 이 계약들에 기대고 있어요(표~\ref{tab:holdout-receiving}). 트레이더가 몇 안 되는 표본에서, 규칙~B와 함께 등록한 비교는 재분석에서 이런 결과를 찾았어요. 주된 기간에서, 나중에 강제 청산을 피하는 일에 대해 허락이 송금보다 그냥 센 수의 수준에서 앞선다는 결과예요. 다만 등록 분석에서는 이 결과를 찾지 못했어요. 넷째, 라벨을 본 뒤에 계산한 대비에서, 어떤 페이지랭크도 기간 밖에서 자기가 다듬는 원시 차수를 이기지 못했어요. 허락 쪽에서는 고르게 다시 출발하기(균등 재시작)로도, 활동 재시작으로도 이기지 못했어요. 송금 쪽에서는 AWP가 그냥 센 수보다 못했어요. 그리고 고르게 다시 출발하는 송금 걷기는 점 추정값(가장 그럴듯한 한 값)으로 보면 그보다도 더 낮았어요. 다시 출발 규칙은 걷기가 자기 화살표를 그냥 센 수에 얼마나 모자라는지만 바꾸었어요. 다섯째, 한 걸음 깊이에서, 고르게 다시 출발하는 걷기의 점수는 상수 하나를 제쳐 두면 스펜더로 들어오는 이동 몫(그 스펜더로 가는 이동 확률의 합)의 일차 함수(곧은 선 모양의 식)예요. 이것이 허락 쪽 결과의 많은 부분을 설명해요(\ref{sub:theory-closed-form}절). 이 주장들을 시험할 수 있게 해 주는 구성 타당도(점수가 재려는 것을 정말 재는지)의 틀은 기여의 한 부분이에요 \parencite{cronbach1955construct,messick1995validity}. 공개된 처리 과정(파이프라인)도 기여의 한 부분이에요. 이 처리 과정이 공개한 요약으로 다른 사람들이 모든 표를 확인할 수 있어요. 탈중앙 시스템이 신용 점수를 잘못 계산하는지는 모두가 관심을 가져야 할 문제이기 때문이에요 \parencite{packin2024credit}.

이 연구의 범위는 일부러 좁게 잡았어요. 이 논문은 온체인 평판이 담보를 대신할 수 있다는 것을 보여 주지 않아요 \parencite{schar2021defi,gudgeon2020crisis}. 또 결과가, 살펴본 기간에 아비트럼 원에서 다섯 토큰으로 허락을 받은 계약을 넘어서까지 들어맞는다는 것도 보여 주지 않아요. 그리고 어떤 점수가 신용도(빌린 돈을 갚을 수 있는 정도), 채무 불이행(빌린 돈을 갚지 못함) 위험, 거래 솜씨를 잰다는 것도 보여 주지 않아요. 또 C-PR이나 엔도스랭크가 시빌 공격(가짜 주소를 많이 만들어 속이는 공격)을 버텨 낸다고 주장하지도 않아요. 고르게 순간 이동(아무 점에서나 다시 출발)하는 걷기에서는, 공격자 주소들로 이루어진 닫힌 농장(공격자가 만든 가짜 주소 무리)이 적은 비용으로 두 점수 어느 것이든 부풀릴 수 있어요(\ref{sec:sybil-cost-model}절).

\dissertationSubheading[sec:future-work]{앞으로 할 일}

아래의 제안은 하나하나가 \ref{ch:discussion}장의 한계에 대한 답이에요. 제안에서 나오는 새 점수나 무게 매기기는, 그것으로 점수를 매기기 전에 정해 두어야 해요. 그리고 등록한 새 자료에서 다시 해 보고(재현하고), 기간 밖에서 차수 기준선(비교 기준)과 견주어 점검해야 해요.

\begin{itemize}
\item 결합 걷기를 등록해서 하는 시험. 2026년 10월부터 12월까지처럼 더 나중 기간의 로그(계약이 남기는 기록)를 뽑기 전에 등록 문서를 커밋(저장소에 확정해 기록함)할 수 있어요. 그런 등록 문서는 수정 분석의 C-PR을, 그 달러 금액, 다섯 토큰, 금액 기울기, 허락 상한(이보다 큰 금액은 이 값으로 자름), 열두 달 점수 기간과 함께 정해 둘 거예요. 그리고 결합 걷기 하나를 더 정해 둘 거예요. 이 걷기의 층은 AWP와 엔도스랭크의 화살표와 무게를 쓰고, 다시 출발할 때는 보내는 활동과 허락하는 활동을 섞어요. 이 걷기가 C-AWP예요. C-AWP의 두 끝(한 층만 걷는 경우)은 AWP와, 활동 재시작을 쓰는 엔도스랭크예요. 또 동결일까지의 이벤트에만 적용하는 코호트 규칙, 규칙~A와~B, 트레이더에 대한 비교도 정해 둘 거예요. 트레이더에 대한 비교에는 W0과 W1에 있던 것보다 더 많은 계약 트레이더(GMX V2에서 거래하는 계약)가 필요해요. C-PR은 결합 방식, 무게, 다시 출발 규칙이 한꺼번에 AWP와 달라요. C-AWP라면 결합 방식 하나만을 이미 발표된 걷기들과 맞세울 수 있을 거예요. 금액 기울기 점검은 홀드아웃 라벨을 다루지 않았어요(표~\ref{tab:robustness-slope}). 그래서 기울기와 상한을 새 기간에 맞추어 조정해서는 안 돼요. 점마다 다른 무게나 학습한 무게, 층마다 다른 감쇠 계수(화살표를 따라갈 확률), 두 정상 벡터(오래 걸었을 때 각 점에 머무는 비율)를 섞은 것은 이 논문이 평가하지 않은 또 다른 결합 방식이에요. 그 가운데 어느 것이든, 그래프 퍼뜨림(전파)을 속을 들여다볼 수 없는 학습 점수와 갈라 주는 검사 가능성(다른 사람이 따져 볼 수 있음)을 지켜야 해요 \parencite{packin2024credit,kotb2023credit}.
\item 더 깊은 표본에서의 퍼뜨림. 여기서 살펴본 에고 네트워크(코호트 계약에 닿는 화살표만 모은 그래프)에서는, 한 걸음 깊이의 스펜더라면 모두 엔도스랭크가 무게를 다시 단 허락 수로 줄어들어요(\ref{sub:theory-closed-form}절). 두 걸음 이상 바깥까지 뽑은 표본으로 시험해 볼 수 있어요. 그 새 자료에서 엔도스랭크와, AWP의 활동 재시작을 쓰는 그 변형을 받은 허락 수와 견주어 시험하는 거예요. 이때 한 걸음 깊이의 스펜더와 더 깊은 스펜더를 따로 시험해요. 그러면 걷기가 무게를 다시 단 수에 그치지 않는 곳에서, 퍼뜨림이 그냥 센 수에 무언가를 보태는지 알 수 있을 거예요.
\item 다른 환경에서 결과 검증하기. 다른 롤업(여러 거래를 묶어 처리하는 체인), 레이어 1 네트워크(가장 바닥에 있는 블록체인), 그리고 더 나중 기간에서 같은 분석을 해 보면, 결과가 아비트럼 원과 살펴본 몇 달을 넘어서도 들어맞는지 알 수 있을 거예요. 디파이는 서비스(프로토콜)마다 모습이 아주 달라요 \parencite{schar2021defi,werner2022sok}. 그래서 이것이 외적 타당도(바깥 자료로도 맞는지)를 위협하는 가장 큰 것이에요.
\item 공격을 가정한 평가. 시빌 비용 모형은 계산으로만 따진 것이고, 체인 위에서 실제로 공격을 해 보지는 않았어요(\ref{sec:sybil-cost-model}절). TrustRank처럼, 순간 이동과 매달린 몫(나가는 화살표가 없는 점에 고인 몫)을 다시 나누는 일을 미리 검증한 씨앗 집합에만 하도록 제한할 수 있어요 \parencite{gyongyi2004trustrank}. 그러면 닫힌 농장이 이용하는, 고르게 흘러드는 몫을 없앨 수 있을 거예요. 그리고 자전 허락(자기끼리 허락을 주고받아 많은 척하기), 짜고 움직이는 스펜더 농장, 전략적인 허락 취소로 레드팀 평가(일부러 공격하는 역할을 맡아 시험하는 평가)를 할 수 있어요. 그러면 공격에 실제로 드는 비용을 재고, 시빌 보정 대리 지표가 여전히 쓸모 있는 정보를 주는지 시험할 수 있을 거예요.
\item 증명 도구와 합치기. C-PR이나 엔도스랭크가 계약에 순서를 매기기 전에, 허락하는 주인들을 패스포트(Passport) 같은 관문으로 걸러 내는 것은 그럴듯한 쓰임새예요 \parencite{humanpassportnddocs}. 이 방법은 시빌인지 아닌지 라벨을 붙인 자료에서 평가해야 해요. 이때 관문 자체의 약점도 함께 살펴야 해요 \parencite{siddarth2020watchmen,bordeianu2026sybil}.
\item 실제로 써 보는 연구. 점수를 위험 오라클(블록체인 밖의 정보를 블록체인에 알려 주는 장치)과 거버넌스(서비스의 규칙을 함께 정하는 일) 도구에 이어 보면, 점수의 실행 시간 특성이 실제 서비스에서 점수를 새로 고치는 주기를 버틸 수 있는지 시험할 수 있을 거예요. 수수료 등급이나 감시 주기를 두고 통제 실험(조건을 맞춘 실험)을 하면, 증거가 \ref{sub:deploy-industry}절의 도입 가설을 대신하게 할 수 있을 거예요.
\end{itemize}

\dissertationSubheading[sec:closing-reflection]{맺으며}

공격자가 있는, 가명의(진짜 이름 대신 주소로만 활동하는) 시장에서, 온체인 평판이 담보의 필요를 없애 주지는 않아요 \parencite{gudgeon2020crisis}. 하지만 온체인 평판은 행동의 구조를 읽을 수 있게 해 줄 수 있어요. 그리고 어떤 구조를 읽는지는 쓰는 화살표에 달려 있어요. 송금은 돈의 흐름을 기록하고, 허락은 맡긴 권한을 기록해요. 페이지랭크는 어느 그래프든 줄 세울 수 있고, 한 걷기가 둘 다 따라갈 수도 있어요. 이 자료에서 결합 걷기는 나중의 스펜더 라벨을 모두 허락만 걷는 걷기보다 더 잘 줄 세웠어요. 하지만 미리 정한 규칙 아래에서는 어느 층도 이기지 못했어요. 그리고 송금만 걷는 걷기보다 나아진 것은 주로 그 걷기가 줄 세우지 못하는 계약들 사이에서였어요. 주인들이 어떤 계약에 권한을 맡겼는지 알려 주는, 빠르고 따져 볼 수 있는 신호가 필요한 서비스(프로토콜)나 연구자가 있을 수 있어요. 그런 쪽은 여기서 살펴본 것과 같은 자료에서 허락 화살표를 쓰면 돼요. 무엇보다 송금 화살표가 없는 스펜더에서 그래요. 허락한 주인 수는 허락 화살표의 가장 값싼 형태예요. 그리고 사후에 보면, 이 수는 두 기간 모두에서 이 연구의 어떤 페이지랭크보다도 나중의 허락을 더 잘 줄 세웠어요. 엔도스랭크는 퍼뜨린 형태이고, 닫힌 농장이 적은 비용으로 부풀릴 수 있어요(\ref{sec:sybil-cost-model}절). 송금이 모이는 중심(허브)을 그려 내고 그곳으로 들어올 흐름을 내다보는 데에는 송금 화살표를 쓰면 돼요. 거기서도, 역시 사후에 보면, 그냥 센 수가 나중에 들어오는 흐름을 AWP보다 더 잘 줄 세웠어요. 강제 청산에 대비해 거래 상대를 가려낼 때는, 주된 기간에서 계약 트레이더 사이에서 허락한 주인 수가 송금 수보다 앞섰어요. 다만 W0에서는 가려지지 않았어요. 그리고 이것을 확실히 하려면 더 많은 트레이더가 필요해요.
